\chapter{Complexity — Random Walk Cover}
%%% generated by the blueprint pipeline from TCSlib.Complexity.Randomized.RandomWalkCover
\section{Overview}

Declarations of the Lean module \texttt{TCSlib.Complexity.Randomized.RandomWalkCover}.

\section{Declarations}

\begin{definition}[Endpoint of a walk]
\lean{Randomized.RegularMultigraph.endpoint}
\label{Randomized.RegularMultigraph.endpoint}
\leanok
Let $G$ be a regular multigraph on a vertex type $V$, with neighbor choices indexed by
$	ext{Fin}(D)$. Starting at a vertex $v$, follow the neighbor indexed by each successive
entry of a finite list of indices; the endpoint is the vertex reached after all choices,
and is $v$ when the list is empty.
\end{definition}

\begin{theorem}[Endpoint of concatenated walks]
\lean{Randomized.RegularMultigraph.endpoint_append}
\label{Randomized.RegularMultigraph.endpoint_append}
\leanok
\uses{Randomized.RegularMultigraph.endpoint}
\difficulty{4}
Let $G$ be a regular multigraph on $V$. For every vertex $v$ and finite lists of
neighbor-slot choices $w,w'$, following the concatenated list $w++w'$ from $v$ ends at
the same vertex as first following $w$ from $v$ and then following $w'$ from that
endpoint.
\end{theorem}

\begin{definition}[Connectivity of a multigraph]
\lean{Randomized.RegularMultigraph.Connected}
\statementsource{Vad12}{pdf-bf2ce5422088-p035-b005}
\label{Randomized.RegularMultigraph.Connected}
\leanok
\uses{Randomized.RegularMultigraph.Adj}
A regular multigraph $G$ is connected if, for every pair of vertices $u,v$, there is a
finite sequence of adjacency steps from $u$ to $v$, allowing a sequence of zero steps
when $u=v$.
\end{definition}

\begin{definition}[Bipartiteness of a multigraph]
\lean{Randomized.RegularMultigraph.IsBipartite}
\label{Randomized.RegularMultigraph.IsBipartite}
\leanok
A regular multigraph $G$ is bipartite if its vertices can be colored with two colors so
that, for every vertex $v$ and every neighbor index $i$, the vertex reached from $v$
using index $i$ has the other color.
\end{definition}

\begin{definition}[Adjacency in the bipartite cover]
\lean{Randomized.RegularMultigraph.coverAdj}
\label{Randomized.RegularMultigraph.coverAdj}
\leanok
Let $G$ be a regular multigraph. Two pairs $(x,c)$ and $(y,d)$, consisting of vertices
and Boolean colors, are adjacent in its bipartite cover exactly when some neighbor index
$i$ satisfies $G$'s neighbor at $x$ indexed by $i$ equals $y$, and $d$ is the opposite
of $c$.
\end{definition}

\begin{theorem}[Symmetry of lifted adjacency]
\lean{Randomized.RegularMultigraph.coverAdj_symm}
\label{Randomized.RegularMultigraph.coverAdj_symm}
\leanok
\uses{Randomized.RegularMultigraph.adj_symm, Randomized.RegularMultigraph.coverAdj}
\difficulty{4}
Let $G$ be a regular multigraph on a vertex set $V$, and give each vertex a Boolean
label. If $(x,b)$ is joined to $(y,c)$ in the two-sheeted lift—meaning some neighbor
slot at $x$ leads to $y$ and $c$ is the complement of $b$—then $(y,c)$ is joined to
$(x,b)$ in the same sense.
\end{theorem}

\begin{theorem}[Lifting a walk to the two-sheeted graph]
\lean{Randomized.RegularMultigraph.cover_lift}
\label{Randomized.RegularMultigraph.cover_lift}
\leanok
\uses{Randomized.RegularMultigraph.Adj, Randomized.RegularMultigraph.coverAdj}
\difficulty{5}
Let $G$ be a regular multigraph on $V$, and suppose there is a finite walk in $G$ from
$v$ to $w$. There is a Boolean value $\pi$ such that, for every starting sheet label
$b$, the lifted walk in the two-sheeted graph goes from $(v,b)$ to
$(w,b\mathbin{\mathrm{xor}}\pi)$.
\end{theorem}

\begin{theorem}[Reachability in the lifted graph]
\lean{Randomized.RegularMultigraph.cover_reach}
\label{Randomized.RegularMultigraph.cover_reach}
\leanok
\uses{Randomized.RegularMultigraph.Connected, Randomized.RegularMultigraph.IsBipartite, Randomized.RegularMultigraph.coverAdj, Randomized.RegularMultigraph.cover_lift}
\difficulty{6}
Let $G$ be a regular multigraph on $V$. Suppose every pair of vertices is connected by a
finite sequence of neighbor steps, and suppose there is no Boolean coloring of $V$ that
assigns opposite colors across every neighbor step. Then for every $u,v	o V$, there is a
finite walk in the two-sheeted lift from $(u,	ext{false})$ to $(v,	ext{false})$, where
each lifted step follows a neighbor step and flips the Boolean label.
\end{theorem}

\begin{theorem}[Words Representing Walks]
\lean{Randomized.RegularMultigraph.word_of_walk}
\label{Randomized.RegularMultigraph.word_of_walk}
\leanok
\uses{Randomized.RegularMultigraph.coverAdj, Randomized.RegularMultigraph.coverAdj_symm, Randomized.RegularMultigraph.endpoint}
\difficulty{5}
Let $G$ be a regular multigraph of degree $D$, and form a graph whose vertices are pairs
$(v,b)$ with $v$ a vertex of $G$ and $b	ext{ is a bit}$, joining $(v,b)$ to $(v',b')$
when some neighbor index $i$ satisfies $v'=	ext{nbr}(v,i)$ and $b'$ is the opposite bit.
For every walk $p$ in this graph from $(x_1,x_2)$ to $(y_1,y_2)$, there is a word $w$ of
neighbor indices of length equal to the length of $p$, whose walk in $G$ goes from $x_1$
to $y_1$, and $y_2$ is $x_2$ toggled if and only if the length of $p$ is odd.
\end{theorem}

\begin{theorem}[An even-length walk]
\lean{Randomized.RegularMultigraph.exists_even_word}
\label{Randomized.RegularMultigraph.exists_even_word}
\leanok
\uses{Randomized.RegularMultigraph.Connected, Randomized.RegularMultigraph.IsBipartite, Randomized.RegularMultigraph.coverAdj, Randomized.RegularMultigraph.cover_reach, Randomized.RegularMultigraph.endpoint, Randomized.RegularMultigraph.word_of_walk}
\difficulty{5}
Let $G$ be a regular multigraph on a finite vertex set $V$. Suppose every pair of
vertices is connected by a finite sequence of neighbor steps, and no Boolean coloring
assigns opposite colors across every neighbor step. For any vertices $u,v$, there is a
list of neighbor-slot choices of even length strictly less than $2|V|$ whose walk from
$u$ ends at $v$.
\end{theorem}

\begin{theorem}[Lengthening a walk by even steps]
\lean{Randomized.RegularMultigraph.exists_word_add}
\label{Randomized.RegularMultigraph.exists_word_add}
\leanok
\uses{Randomized.RegularMultigraph.adj_symm, Randomized.RegularMultigraph.endpoint, Randomized.RegularMultigraph.endpoint_append}
\difficulty{4}
Let $G$ be a regular multigraph of positive degree, let $u,v$ be vertices, and let $w$
be a finite sequence of neighbor indices whose walk from $u$ ends at $v$. For every
natural number $m$, there is another sequence of neighbor indices whose walk from $u$
also ends at $v$ and whose length is the length of $w$ plus $2m$.
\end{theorem}

\begin{theorem}[A walk of prescribed length]
\lean{Randomized.RegularMultigraph.exists_endpoint_eq_two_mul}
\label{Randomized.RegularMultigraph.exists_endpoint_eq_two_mul}
\leanok
\uses{Randomized.RegularMultigraph.Connected, Randomized.RegularMultigraph.IsBipartite, Randomized.RegularMultigraph.endpoint, Randomized.RegularMultigraph.exists_even_word, Randomized.RegularMultigraph.exists_word_add}
\difficulty{2}
Let $G$ be a regular multigraph on a finite, nonempty vertex set $V$, with positive
degree. Suppose every pair of vertices is connected by a finite sequence of neighbor
steps, and no Boolean coloring assigns opposite colors across every neighbor step. For
any vertices $u,v$, there is a list of neighbor-slot choices of length exactly $2|V|$
whose walk from $u$ ends at $v$.
\end{theorem}
